\documentclass{article}
\usepackage[centering,tmargin=30mm,bmargin=20mm,lmargin=1.0in,rmargin=1.0in]{geometry}
\usepackage[hidelinks]{hyperref}
\usepackage{wallpaper}
\usepackage{amsmath}
\renewcommand{\familydefault}{\sfdefault}

\date{\today}
\makeatletter
\let\thedate\@date
\makeatother

\setlength\parindent{0mm}
\setlength\parskip{2mm}

\begin{document}

\thispagestyle{empty}
Managing Editor\\
\textit{Nano Letters}\\
United States\\

\vspace{5mm}

Dear Prof Besley, 

We thank you and the reviewers for providing this feedback on our submitted manuscript. We have made many significant revisions, 
and believe that the manuscript is now in much better shape. Please find a detailed response to each reviewer's comments and a summary 
of all changes made below.


Reviewer 1:

% Comment 1:
\textbf{Comment}

1.	The first significant concern pertains to the manuscript's title, 
which suggests a general solution for molecular counting. 
However, the methodology appears to be specifically tailored for Fluorescence Correlation Spectroscopy 
(FCS) with a rigid two-state Markov Chain structure. 
This narrow focus may mislead readers regarding the method's applicability to other SMLM 
techniques such as PALM and dSTORM, where a two-state model has proven inadequate. 
It is recommended that the authors revise the title to more accurately reflect the scope of 
their methodology, emphasizing its application to FCS or clarifying its limitations concerning other SMLM techniques.

\underline{\textbf{Response:}}
We agree with the reviewer that the manuscripts title may be too general and as a result misleading for some readers.
To clarify this, we have changed the title to: "A Bayesian model to count the number of two-state emitters in a 
diffraction limited spot". 

Changes:
\begin{description}
    \item[$\bullet$] Changed the title to: A Bayesian model to count the number of two-state emitters in a diffraction limited spot
\end{description}

% Comment 2:
\textbf{Comment}

2. Second, and in my opinion, the introduction of the manuscript inadequately addresses the 
breadth of existing quantitative methods for molecular counting in SMLM. Notably, significant works 
(including but not limited to Rollins et al. (2014), Patel et al. (2021) and Nino et al. (2017, 2019 2021)) 
are omitted, which neglects important contextual frameworks and benchmarks against which 
blinx's performance should be compared in the more general scope of SMLM. 
Expanding the literature review to include a discussion of these existing methods and providing 
comparative analyses would better highlight blinx's advantages and limitations. 
Furthermore, the manuscript incorrectly claims that prior methods offer only single estimates of 
molecular counts, disregarding current approaches which indeed utilize Bayesian estimates. 
While this statement is factually inaccurate relative to the title of the manuscript, if the 
focus is on counting in FCS, this should be made much clearer and the introduction should indeed 
reflect this the positioning of blinx as a unique probabilistic estimator for FCS.

\underline{\textbf{Response:}}
We thank the reviewer for highlighting these highly relevant references. We have re-written a section of 
the introduction to include these references and highlight their previous work in Bayesian modeling of 
these systems and in Markovian modeling of blinking fluorophores.

We also thank the reviewer for raising a concern over the claim of other methods. We have changed the text to 
clarify that the other DNA-PAINT based models, that are directly comparable to blinx, offer only a single estimate 
of count. As the reviewer correctly pointed out, many of these other Bayesian methods do provide a probabilistic 
estimate.

Changes
\begin{description}
    \item[$\bullet$] Changed "all prior methods provide a single estimate of count" to 
    "these DNA-PAINT based methods provide a single estimate of molecular count"
    \item[$\bullet$] Re-wrote the discussion of other Bayesian methods in the introduction 
    to highlight specific contributions from each of the suggested references. 
    \item[$\bullet$] Added references [31, 34, 35, 36]
\end{description}

% Comment 3
\textbf{Comment}

3.	On the topic of emphasizing the novelty of Bayesian modeling, 
the third critical issue arises from the manuscript's assumption of a uniform (flat) 
prior on the number of emitters within a Bayesian framework. 
This effectively reverts to a frequentist (i.e.\ non-Bayesian) estimate, 
negating the benefits of Bayesian methods such as incorporating prior knowledge 
and handling low-probability scenarios which may arise in FCS or otherwise, 
but most importantly contradicting the supposed novelty of Bayesian modeling for this application. 
The authors should reconsider the choice of a uniform prior and either explore or discuss 
the integration of other informative priors that leverage existing knowledge or distributional 
assumptions to fully utilize the strengths of the Bayesian framework.

\underline{\textbf{Response:}}
We failed to communicate the role of the prior clearly: a uniform or flat prior on the number of emitters is 
not an assumption of our model. In the original manuscript we imprecisely said that this prior is uniform and 
can be ignored. But what we meant to convey is that this prior is application specific and can be set arbitrarily. For our experiments (where we do not wish to impose a prior on $n$ for a fair evaluation of our model), we use a uniform prior. We have updated the text to make 
this point more clear. 

Concretely, we used a uniform prior between 1 and 35 emitters for the simulated experiments, and between 1 and 8 for the DNA origami experiments.
We changed the text of the results section to include this information.
We chose a uniform prior because we wanted to demonstrate the robustness of our model, but one could easily 
add more informative priors to the model, without fundamental changes to the derivation shown in the methods. 

Changes:
\begin{description}
    \item[$\bullet$] Stated that $p(n)$ is application specific in section 2.1 Model
    \item[$\bullet$] Explained in section 3.1 Simulated Experiments that we placed a uniform prior between 1 and 35
    \item[$\bullet$] Explained in section 3.4 Experimental Counting that we placed a uniform prior between 1 and 8
\end{description}

% Comment
\textbf{Comment}

 4.	Fourth, the necessity of employing the Laplace approximation within the 
 Bayesian inference framework is questionable, particularly given the presence of a flat prior. 
 Maximum Likelihood Estimation (MLE) with Newton iterations could potentially suffice without 
 introducing the additional complexity of the Laplace approximation. 
 The authors should provide a rationale for choosing the Laplace approximation over 
 alternative inference methods, particularly since many approximations to the true 
 posterior are made in model construction. If the approximation does not offer clear advantages, 
 simplifying the inference process by adopting standard MLE techniques should be considered.

 \underline{\textbf{Response:}}
 We would like to clarify that we are using Maximum Likelihood Estimation, and optimizing our model 
 by performing gradient ascent directly on our loss function. We employ the Laplace Approximation 
 post-hoc to normalize likelihoods between different models (different values of $n$): because we model each molecular count 
 individually (with a different number of possible states) the model likelihoods are not comparable 
 in principle. The Occam's factor normalizes models of differing complexity, allowing us to directly 
 compare the likelihoods of different models and construct the posterior.

 In order to clarify this we now explicitly state in the text that we first find the MLE, and then use Occam's factor to compensate 
 for different model complexities.

 Changes
 \begin{description}
    \item[$\bullet$] Re-structured section 2.1.3 Inference
\end{description}


% Comment 5:
\textbf{Comment}

5.	Fifth, the model construction lacks adequate references to foundational concepts 
like Poisson noise or photon intensity distributions via Point Spread Functions (PSFs). 
This omission limits the reader's ability to understand the theoretical underpinnings 
and validate the model's assumptions. Incorporating additional references would provide a 
more comprehensive background on the noise models and intensity distributions utilized in blinx's construction.

\underline{\textbf{Response:}} 
We have added references to both Poisson and Gaussian noise when describing the
model construction. To further increase clarity we have added a few sentences
that describe how we obtain our intensity measurement $x$ from the image. 

Changes
\begin{description}
    \item[$\bullet$] Added references [37, 38, 39] to section 2.1.1 Intensity Model. 
    Discussing the contributions of Poisson noise and Gaussian noise to the system.
    \item[$\bullet$] Expanded section 2.1.1 Intensity Model explaining 
    how traces were extracted from the raw images.  
\end{description}

% Comment 6
\textbf{Comment}

6.	Sixth, the validity of using a large-sample Gaussian approximation for the 
Poisson-distributed shot noise is questionable, especially for very short frame 
intervals where photon counts are low. In scenarios with extremely high frame 
rates and minimal photon counts per frame, this approximation becomes invalid, 
potentially affecting downstream analysis accuracy. The authors should assess the 
range of frame intervals and photon counts where the Gaussian approximation remains valid. 
If significant deviations occur, alternative approaches should be considered, or the 
limitations of the approximation should be explicitly discussed. 

\underline{\textbf{Response:}}
We believe that the Gaussian approximation is very reasonable especially in the regime where our model operates. 
To demonstrate this point we plotted the KL-divergence between the Poisson and Gaussian distributions 
as a function of photon count $\lambda$ (see Figure S1). The absolute KL-divergence at which the approximation fails 
is debatable, but we chose a value of $10^{-3}$ which corresponds to a $\lambda$ value of 81. We experimentally 
observed $\lambda=2084$ which corresponds to a KL-divergence less than $10^{-4}$.
%
Based on our calculated emission rates, a frame rate greater than 250 frames
per second would be required to drop below a photon count of 81. At such high
frame rates we expect that low SNR would become the limiting factor rather than
this approximation. Also note that this is the same approximation used in Huang
et.\ al.\ 2013.

Changes
\begin{description}
    \item[$\bullet$] Added Figure S1 and referenced in section 2.1.1 illustrating the KL-divergence between 
    the Poisson and Gaussian distributions as a function of photon count ($\lambda$)
    \item[$\bullet$] Added section S9 to the appendix and referenced in section 2.1.1 (Intensity Model discussion) the validity of this 
    approximation
    \item[$\bullet$] Referenced Huang et. al. 2013 [38] in section 2.1.1 Intensity Model, where the same 
    approximation is used
\end{description}

% Comment 7
\textbf{Comment}

7.	Seventh, the manuscript assumes time as discrete and posits that all 
emitters share identical and constant blinking parameters (pon and poff). 
This is unrealistic for general imaging techniques, where pon is much less than poff is 
typical (and calibrated) to ensure sparse localizations for state-of-the-art 
SMLM to work as effectively as possible. These strong assumptions may not hold 
in practical FCS applications, leading to inaccurate molecular counts and limited 
generalizability (as is clearly shown in the later studies of the proposed model) 
to more diverse experimental conditions.

\underline{\textbf{Response:}}
We would like to clarify on what is meant by "identical and constant blinking parameters". 
Our model assumes that the $p_\text{on}$ of an emitter in a spot is equal to that of every other 
emitter in the same spot, and that the same is true for $p_\text{off}$. However, we do not assume 
that $p_\text{on} = p_\text{off}$ for any emitters. We also address this assumption in the discussion 
clarifying that in real systems this may not be the case, and that we expect a decrease in model 
performance as the parameters of emitters within a single spot diverge. 

In Figure 3a, we evaluate the models performance for a wide variety of $p_\text{on}$ and $p_\text{off}$ 
combinations. Further, in section 3.3 (qPAINT Kinetic Regime) and Figure 4, we evaluate the model 
in the exact regime that the reviewer describes, where $p_\text{on}$ is much less than $p_\text{off}$. 
We find that while the base model fails in this regime, performance is recovered, and even 
surpasses that of qPAINT, once priors are placed on the blinking parameters. 

% Comment
\textbf{Comment}

8.	Due to the nature of the aforementioned approximations, in the simulation analysis, 
noticeable discrepancies exist between the inferred model and raw intensity histograms. 
Additionally, figures, particularly Fig. 2d-e, lack necessary labels and descriptions, 
obscuring the evaluation of the model's performance. Enhancing the description and 
clarity of simulation results, ensuring all figures are fully labeled, and providing 
comprehensive captions would facilitate better understanding. 
The authors should address the apparent discrepancies between the model and simulations, 
possibly by refining the model, quantifying the approximation errors and/or offering 
explanations for the observed differences under a range of parameters choices.

\underline{\textbf{Response:}}
We are not sure what discrepancies between the inferred model and raw intensity histograms 
the reviewer is referring to. We show intensity histograms (in purple) and fit model 
distributions (in orange) in Figures 2a, 2b, 3c, 3d, 4a, 4b, 5a, and 5b, and they closely 
match in each case. For the experimental data shown in Figures 5a and 5b, the intensity histogram 
is slightly noisier, but our inferred model still accurately fits the data. 

In order to clarify each of the figures, especially Figures 2d and 2e, we have added 
labels and descriptions, as well as expanded each of the captions.

Changes
\begin{description}
    \item[$\bullet$] Figure 1: Labeled the posterior distribution and expanded the caption
    to describe that each spot is fit independently
    \item[$\bullet$] Figures 2a, 2b, 3c, 3d, 4a, 4b, 5a, and 5b: the inset panel on the 
    far right has been labeled "posterior" and the y-axis has been added.
    \item[$\bullet$] Figure 2d and 2e: the y axis has been labeled as "estimated n"
    \item[$\bullet$] Figure 2d: additional x-ticks have been added to clarify the trace 
    lengths tested
    \item[$\bullet$] Figure 2: Caption expanded to explain the colorbar in Figure 2c and 
    that Figure 2c shows an averaged posterior over 10 traces for each n.
    \item[$\bullet$] Figure 3: Caption expanded to explain that blinx fails for the kinetic 
    conditions shown in Figure 3c, but works well for those shown in Figure 3d.
    \item[$\bullet$] Figure 4: Caption expanded to explain that this regime poses a challenge 
    for blinx because of the limited number of intensity states observed.
    \item[$\bullet$] Figure 5: Caption expanded to reference Figure S2, and provide a deeper 
    look at the posterior shown in Figure 5e
    \item[$\bullet$] Figure 5b, 5d, 5e: the x-axis of the posterior plots was extended to 
    show the range of the prior placed on possible count $p(n)$.
\end{description}

% Comment
\textbf{Comment}

9.	The manuscript does not account for the effects of adding a noise floor (low SNR) 
which could mitigate inaccuracies in molecular count predictions under low SNR conditions, 
as is also clearly shown in simulations. Ignoring SNR variations may limit the model's 
robustness and applicability in experiments with varying noise levels. Integrating SNR 
considerations into the model, potentially by adapting noise handling mechanisms from 
existing methods, would improve molecular count predictions in low SNR scenarios.

\underline{\textbf{Response:}}

We thank the reviewer for raising this suggestion. In our intensity model we model noise along every step in the 
lightpath of the experiment (readout noise in the detector, shot noise of the fluorophore, and
background photon emission from various sources). However, in experimental traces we still occasionally see some 
unaccounted for noise. To account for these observations, our implementation 
of the noise model includes a minimum probability of observing any intensity 
at any time. This prevents a small fraction of observed intensities that are out-of-distribution 
from having a significant effect on the counting ability of the model. We have added a section on this minimum 
probability to the supplementary information and referenced it in the discussion.

As for experiments with varying noise levels, we can think of varying noise in across regions of the image, 
or noise that varies in time. In terms of noise varying in a single image, our model fits each detected spot 
independently. This allows different regions of the image to have differing noise levels.
%
Our model does not account for noise varying in time, but we believe that with relatively simple extensions it 
it could account for this. However, we believe that this is out of scope for this manuscript and best left 
for any follow-up work. 

Changes
\begin{description}
    \item[$\bullet$] Added section S3.11 Noise Model Inference, to discuss our implementation of a minimum probability for 
    observing any intensity. 
\end{description}

\textbf{Comment}

10. Last, the experimental validation using DNA origami with known counts (N=1 and N=4) 
demonstrates suboptimal posterior probabilities as outputted from blinx, 
indicating the model's inability to assign high confidence to correct counts 
even in simplistic scenarios. For instance, a posterior probability of 0.6 for N=4 
in the real-data analysis (in which it is visually clear that the posterior probability 
that N=4 should be much higher than 0.6), suggests uncertainty that undermines the model's 
reliability, effectiveness and overall utility. The authors should investigate the causes 
of low posterior probabilities and aim to refine the model to enhance its confidence in 
molecular count assignments. Expanding the experimental validation to include a broader 
range of counts and more complex configurations would also provide a more thorough assessment 
of the model's capabilities.

\underline{\textbf{Response:}}
The posterior distribution shown in Figure 4e is the average of the posterior 
calculated for each of the 110 experimental traces. 
For “high confidence” traces, like the one shown in Figure 4b, 
the posterior probability approaches 1, as shown in the inset to the histogram in Figure 4b. 
For other traces, the model is less confident (shown in Figure S2a). 
For these traces the correct count (4) is still the highest probability (~0.7), 
but counts greater than 4 also have a significant probability. 
This behavior is expected. 
One can easily imagine that the true count is 5, but only 4 of the emitters are 
ever “on” at the same time, and thus the highest count ever observed is 4. 
This scenario is less likely than the true count being 4, but is not impossible 
(our model estimates the probability at ~20\%). 
For still other traces the model can be incorrect. 
For example the trace shown in Figure S2b is estimated to have a count of 3 with high probability. 
When averaging across all traces, these incorrect posteriors reduce the average posterior of N=4. 

Changes:
\begin{description}
    \item[$\bullet$] Added Figure S2, to give more context to the distribution shown in Figure 4e. Whereas Figure 4e 
    shows the average of 110 posteriors, this figure shows individual posteriors for 2 example traces.
    \item[$\bullet$] Added section S3.9 "N=4 Experimental Posteriors" to discuss in more detail the posteriors of the N=4 experimental results.
\end{description}

\textbf{Minor Comment}
While the manuscript is generally well-written, certain sections could benefit from clearer 
explanations, particularly in describing the transition and intensity models. Additionally, 
the computational demands of the blinx model are not discussed, which is relevant for practical 
applications involving large datasets.

\underline{\textbf{Response:}}
We agree with the reviewer that discussing the computational demands of blinx is important and
have added a section to the supplemental information, and referenced it in the discussion.

Additionally, we hope that the changes we have made to the methods section in response to other 
comments has increased it's clarity.

Changes:
\begin{description}
    \item[$\bullet$] Added section S3.13 Run Time to discuss the computational demands of blinx
\end{description}

Reviewer 2:

\textbf{Comment}

1. This work should be put in context of the Cox paper 
“Bayesian localization microscopy reveals nanoscale podosome dynamics” 
which takes a similar approach but applies it to the entire image, not just a sub-location.

\underline{\textbf{Response:}}
We thank the reviewer for highlighting this missing reference. We have included it in our discussion of 
other Bayesian approaches to single-molecule imaging in the introduction.

Changes
\begin{description}
    \item[$\bullet$] Added reference and short sentence to the introduction.
\end{description}

\textbf{Comment}

2. The definition of SNR is simple enough to put in the main text and would motivate the discussion, 
which is basically if the intensity peaks in the histogram start to overlap or not. 
Then its intuitively clear why low SNR doesn’t work, because the discrete states in 
intensity go away, and there become large correlations in model parameters to explain the same data. 
It’s also clear why it doesn’t work at low duty cycle, because it’s not enough information to 
distinguish between rates and emitters.

\underline{\textbf{Response:}}
We have moved our definition of SNR to the main text (section 3.1) and removed the associated section from the 
supplemental material.

Changes
\begin{description}
    \item[$\bullet$] Moved definition of SNR to main text
\end{description}

\textbf{Comment}

3. Its worth mentioning  that some assumptions here can be broken. 
The DNA-PAINT docking strands can be damaged over long runs and would appear 
similar to a photobleaching event. 
This can be ameliorated somewhat by buffers, etc and this should be mentioned. 
The assumption of an equal rate of emission for all fluorophores may not be true when imaging real samples. 
For example, when imaging a membrane protein that could be varying distances from the coverslip. 
The emission rate would depend on the TIRF evanescence field and fluorophore distance from the coverslip. 
Even without TIRF, the emission to the coverslip is somewhat depth dependent. 
This is not a consideration for the data shown because its origami flat on the coverslip.

\underline{\textbf{Response:}}
We thank the reviewer for highlighting these possible sources of variation. In the discussion we have expanded our 
section on possible sources of variation in parameters within a spot to include the emission rates. Additionally 
we added a section to the appendix highlighting the potential flaws in the assumption that DNA-PAINT is resistant 
to photo-bleaching. Finally, we would like to point out that we used specific buffers to ameliorate this effect in 
our experimental procedure, and added a reference to that section in the new SI section. 

Changes
\begin{description}
    \item[$\bullet$] Added section S3.10 to the SI discussing limitations of DNA-PAINT, and referenced 
    it from the discussion
    \item[$bullet$] Added reference [49] highlighting this concern with DNA-PAINT
    \item[$\bullet$] Added details about how emission rate could be non-constant to discussion
\end{description}

% \bibliographystyle{unsrt}
% \bibliography{references}

\vspace{5mm}
Thank you, \\
Alex Hillsley,\\
Jan Funke,\\
HHMI Janelia Research Campus

\end{document}
